\section{The planar projection estimate}
\label{sec:projection}

The geometric input is the discrete planar Furstenberg theorem of
Ren and Wang. We state precisely the consequence used below, including
weighted populations, extra parameter labels, and incomplete incidence
graphs. No restriction estimate in three dimensions is used as an input.

For a bounded set \(E\), let \(\abs E_\rho\) denote the number of
half-open dyadic \(\rho\)-squares meeting \(E\), where replacing a scale
by a comparable dyadic scale is harmless. A finite family of dyadic
\(\rho\)-squares \(\mathcal P\) is a \((\rho,d,C)\)-set if
\[
 \#\{p\in\mathcal P:p\cap B(z,r)\ne\varnothing\}
 \le C r^d\#\mathcal P,\qquad \rho\le r\le1,
\]
up to a fixed enlargement of the ball. We use dyadic tubes in a fixed
bounded slope chart. The same definition for tubes refers to their
slope--intercept parameters. For tubes meeting one fixed
\(\rho\)-square, this nonconcentration condition is equivalent, within
fixed factors, to the corresponding condition on slopes.

\begin{theorem}[Ren--Wang, discrete form]
\label{thm:ren-wang}
Fix \(0<s\le1\), \(0<d\le2\), and \(\alpha>0\).
There are \(\eta_0=\eta_0(\alpha,s,d)>0\) and
\(\rho_0=\rho_0(\alpha,s,d)>0\) with the following property.
Suppose \(\mathcal P\) is a nonempty
\((\rho,d,\rho^{-\eta_0})\)-set of squares, with
\(0<\rho<\rho_0\). For each \(p\in\mathcal P\), let
\(\mathcal T(p)\) be a family of tubes meeting \(p\), of cardinality
comparable to \(M\), and forming a
\((\rho,s,\rho^{-\eta_0})\)-set. Then
\begin{equation}
 \#\bigcup_{p\in\mathcal P}\mathcal T(p)
 \gtrsim_\alpha
 \rho^{-\min\{1,d,(d+s)/2\}+\alpha}M .
 \label{eq:ren-wang-discrete}
\end{equation}
The statement is unchanged by fixed bounded domains, finitely many
coordinate charts, or fixed-factor enlargements of the grids and tubes,
after decreasing \(\eta_0,\rho_0\) if necessary.
\end{theorem}

This is Theorem~4.1 of \citet{RenWang2025}, expressed using its
Definition~3.1 of a nice configuration. That definition already permits
each $\mathcal T(p)$ to be a selected subfamily of the tubes meeting $p$.
The reduction below additionally handles probability weights, auxiliary
tags that may share a slope, and incidence graphs of positive product
mass. The fixed-factor variants follow by taking comparable dyadic scales,
covering by boundedly many grid
pieces, and absorbing the resulting fixed constants into the input
loss. The order of the quantifiers is important: the desired output
loss \(\alpha\) is fixed before the admissible input loss \(\eta_0\).

A tag is a slope together with any auxiliary labels; distinct tags
may have the same slope. Keeping the full tag allows its set of
neighboring points to depend on that additional information.

\begin{theorem}[Weighted projection bound with parameter tags]
\label{thm:planar-projection}
Fix \(0<d\le2\), \(0<s\le1\), and \(\alpha>0\), and fixed bounded
point and slope domains. There are \(\eta>0\) and \(\rho_1>0\) such
that the following holds for \(0<\rho<\rho_1\).
Let \(\mu\) be a finite probability population of points in \(\RR^2\).
Let \(\nu\) be a finite probability population of tags \(b\), with a
slope \(a(b)\) assigned to each tag. Suppose
\begin{align}
 \mu(B(z,r))&\le \rho^{-\eta}r^d,
 & (a_*\nu)(I)&\le \rho^{-\eta}r^s
 \label{eq:projection-populations}
\end{align}
for all balls and intervals of radius or length comparable to
\(r\in[\rho,1]\). Let \(G\) be a graph of point--tag incidences with
\((\mu\times\nu)(G)\ge\rho^\eta\).
For each tag \(b\), suppose that the projection
\(\pi_{a(b)}(x)\) of its neighbors can be covered by at most \(M\)
intervals of length \(\rho\). Here
\(\pi_a(x_1,x_2)=x_2-a x_1\). One may multiply these maps by
scalars whose absolute values have fixed positive lower and upper bounds.
Then
\begin{equation}
 M\ge \rho^{-\min\{1,d,(d+s)/2\}+\alpha}.
 \label{eq:weighted-projection-finite}
\end{equation}
Changing the implicit fixed factors in the domains, projections,
and interval lengths changes only \(\eta,\rho_1\).

In particular, along any sequence with subpower nonconcentration
constants, graph product mass \(\rho^{o(1)}\), and conditional
projection count at most \(\rho^{-l-o(1)}\), one has
\begin{equation}
 l\ge \min\{1,d,(d+s)/2\}.
 \label{eq:weighted-projection-limit}
\end{equation}
The limiting assertion for \(d=0\) is immediate.
\end{theorem}

\begin{proof}
Absorb any allowed scalar factors into the interval lengths, and use
\(\pi_a(x_1,x_2)=x_2-a x_1\). Normalized orthogonal projections
also have this form on boundedly many direction charts, after an
orthogonal coordinate change and such a bounded scalar factor. Write
\(q=(\mu\times\nu)(G)\), and independently sample
\[
 N=\lceil\rho^{-d}\rceil\quad\hbox{points},\qquad
 Q=\lceil\rho^{-s}\rceil\quad\hbox{whole tags}.
\]
The tags, including any information besides their slope, are sampled
without alteration.

We give the finite loss accounting. In this proof \(C\) denotes a
fixed numerical constant, enlarged a finite number of times; it does
not depend on \(\rho,\eta\), or the population cardinalities.
For fixed \(\eta>0\) and sufficiently small \(\rho\), the sampled
multisets have, simultaneously for all dyadic balls or intervals at
scales between \(\rho\) and \(1\),
\begin{align}
 \#\{x_i\in B(z,r)\}
 &\le C\rho^{-3\eta}(r/\rho)^d, \label{eq:sampled-point-count}\\
 \#\{a(b_j)\in I\}
 &\le C\rho^{-3\eta}(r/\rho)^s. \label{eq:sampled-slope-count}
\end{align}
Moreover, the failure probability is smaller than every fixed power
of \(\rho\). Indeed each count is binomial, its expectation is bounded
by \(\rho^{-\eta}N r^d\), or \(\rho^{-\eta}Q r^s\), and the displayed
threshold has an additional factor at least a fixed multiple of
\(\rho^{-2\eta}\). The standard exponential Markov estimate
\[
 \PP(X\ge t)\le (e\,\EE X/t)^t,\qquad t\ge \EE X,
\]
follows by applying Markov's inequality to
\(\exp(\theta X)\) and optimizing \(\theta\).
Here the smallest nonvacuous threshold is a positive multiple of
\(\rho^{-3\eta}\). A union bound over the polynomially many dyadic
cells and scales proves the assertion. Every ball or interval is
covered by boundedly many such cells of comparable scale.

The expected sampled edge count is \(qNQ\), whereas its maximum is
\(NQ\). The occupancy failure event therefore contributes at most
\(o(qNQ)\) to this expectation. Some sample obeys
\eqref{eq:sampled-point-count}--\eqref{eq:sampled-slope-count} and
has at least \(qNQ/2\) edges. Delete point vertices of degree
less than \(qQ/4\). At least \(qNQ/4\) edges remain, hence there
are at least \(qN/4\) remaining point vertices.

A finest point square contains at most \(C\rho^{-3\eta}\)
sampled vertices. Keep one remaining vertex in each occupied square;
the resulting point-square family has cardinality at least
\(C^{-1}q\rho^{3\eta}N\). Each selected vertex still has at least
\(qQ/4\) incident sampled tags.

For a tag and one of its allowed projection intervals, its slope
and projection bin specify a tube through the corresponding incident
point. Rounding these data to the tube grid and using a bounded
enlargement costs only fixed factors. At one selected point square,
only boundedly many different intercept bins in a fixed slope bin
can give tubes through that square. By
\eqref{eq:sampled-slope-count}, a tube there can be represented
by at most \(C\rho^{-3\eta}\) of the sampled tags, up to the same
bounded neighboring-bin multiplicity. Remove such coincidences.
Each selected point is incident to at least
\(C^{-1}q\rho^{3\eta}Q\) distinct tubes.

Pigeonhole the point vertices according to dyadic ranges of this
last degree. There are \(O(\log(1/\rho))\) ranges. Keep a range
containing at least that fraction of the selected points, and prune
the tube families to a common size \(M_0\) within a factor of two.
For small \(\rho\), all these operations give
\begin{equation}
 \#\mathcal P\ge \rho^{C\eta}N,\qquad
 M_0\ge \rho^{C\eta}Q .
 \label{eq:projection-retained-sizes}
\end{equation}
The original sampled count bounds remain upper bounds on every
retained subset. Dividing them by the retained cardinalities in
\eqref{eq:projection-retained-sizes} proves that \(\mathcal P\)
is a \((\rho,d,\rho^{-C\eta})\)-set and every
\(\mathcal T(p)\) is a
\((\rho,s,\rho^{-C\eta})\)-set. The constants can be taken uniform
through the finite grid changes.

On the other hand, each of the \(Q\) sampled tags has at most \(M\)
allowed projection bins, so the total number of distinct tubes is
at most \(CQM\). Put
\[
 \chi=\min\{1,d,(d+s)/2\}.
\]
Apply Theorem~\ref{thm:ren-wang} with output loss \(\alpha/4\).
Choose \(\eta\) sufficiently small that
\(C\eta<\eta_0(\alpha/4,s,d)\) and \(C\eta<\alpha/4\), and
then choose \(\rho_1\) sufficiently small for all preceding
bounds and fixed constants. The theorem yields
\[
 CQM\gtrsim_{\alpha}
 \rho^{-\chi+\alpha/4}M_0
 \ge \rho^{-\chi+\alpha/4+C\eta}Q .
\]
Absorbing the fixed constants by another loss less than \(\alpha/2\)
gives \eqref{eq:weighted-projection-finite}.

For the last assertion, apply the finite theorem at any fixed
\(\alpha>0\). Its hypotheses eventually hold because their losses
are subpower. Comparison with the assumed projection count gives
\(l\ge\chi-\alpha\); let \(\alpha\downarrow0\).
\end{proof}

\begin{remark}[Incidence laws and conditional applications]
\label{rem:projection-incidence-laws}
Often an incidence graph is supplied by a trimmed subprobability
\(\sigma\) rather than directly by \(\mu\times\nu\).
If \(\sigma(G)\ge m\) and \(\sigma\le K\mu\times\nu\), then
\[
 (\mu\times\nu)(G)\ge m/K.
\]
Thus subpower mass and subpower product domination give exactly
the graph hypothesis of Theorem~\ref{thm:planar-projection}.
Uniform lower bounds on retained conditional atom masses give the
required projection support counts. Uniform upper bounds at every
fixed prefix, followed by a sufficiently fine finite scale mesh,
give the required ball-mass bounds. The conditioning and
homogenization results in Section~\ref{sec:entropy} justify those
operations when applied in the prescribed order. In each application
below we specify which point population and parameter population are
being compared; an entropy lower bound without this uniformization
is not a substitute for \eqref{eq:projection-populations}.
\end{remark}
